\section[Introduction]{Introduction}

\begin{frame}{Contenu du cours}
	L'analyse en composantes principales (ACP) est un algorithme d'\textbf{apprentissage non supervisé} qui a pour objectif de réduire la \textbf{dimensionnalité} d'un problème de régression ou de classification. Dans ce contexte, la dimensionnalité est liée aux caractéristiques d'entrée $x \in \mathcal{X}$.
	\vspace{0.5cm}
	
	Dans ce cours, nous aborderons les éléments suivants:
	\begin{itemize}
		\pause
		\item Optimisation sous contraintes avec les \textbf{multiplicateurs de Lagrange}.
		\pause
		\item Problématique de la réduction de dimensions pour l'\textbf{apprentissage statistique}.
		\pause
		\item Propriétés de l'\textbf{algèbre linéaire} utiles pour l'ACP.
		\pause
		\item Application pratique au \textbf{traitement d'images}.
	\end{itemize}
\end{frame}
